% UC-4 Ordinamenti — Template di Cockburn (LONGTABLE con griglia completa)
% Include con: % UC-4 Ordinamenti — Template di Cockburn (LONGTABLE con griglia completa)
% Include con: % UC-4 Ordinamenti — Template di Cockburn (LONGTABLE con griglia completa)
% Include con: % UC-4 Ordinamenti — Template di Cockburn (LONGTABLE con griglia completa)
% Include con: \input{UC3 - Cockburn.txt}
% Compilazione: XeLaTeX o LuaLaTeX consigliati (oppure pdfLaTeX con \usepackage[utf8]{inputenc})
% Preambolo richiesto:
% \usepackage{longtable}
% \usepackage{array}
% \usepackage{enumitem} % per itemize/enumerate compatti
\renewcommand{\arraystretch}{1.25}
\setlength{\tabcolsep}{8pt}

\begin{longtable}{|p{4cm}|p{11cm}|}
\caption{UC-3 Ordinamenti -- Template di Cockburn}\label{tab:uc3-cockburn}\\
\hline
\textbf{Campo} & \textbf{Valore} \\ \hline
\endfirsthead

\hline
\textbf{Campo} & \textbf{Valore} \\ \hline
\endhead
\hline
\multicolumn{2}{r}{\small Continua nella pagina successiva} \\
\endfoot
\hline
\endlastfoot

\textbf{ID - Nome dello UC} & 3 - Ordinamento \cite{uc3-figma} \\ \hline
\textbf{Data di creazione} & 14/10/2025 \\ \hline
\textbf{Attore primario} & Utente \\ \hline
\textbf{Trigger} & L'utente indica di voler ordinare l'elenco delle issues. \\ \hline
\textbf{Descrizione} & L'utente vuole ordinare le issues presenti nell'elenco tramite un criterio specifico. \\ \hline

\textbf{Precondizioni} &
\begin{minipage}[t]{\linewidth}\raggedright
\begin{itemize}[leftmargin=*,nosep]
\item PRE-1. L'utente è stato registrato da un amministratore ed ha effettuato correttamente l'autenticazione.
\item PRE-2. Il sistema ha caricato la pagina di riepilogo delle issues.
\end{itemize}
\end{minipage} \\ \hline

\textbf{Postcondizioni} & L'elenco delle issues è riordinato. \\ \hline

\textbf{Scenario di successo} &
\begin{minipage}[t]{\linewidth}\raggedright
\begin{enumerate}[leftmargin=*,nosep]
\item L'utente visualizza la tabella contenente le issues con le relative colonne.
\item L'utente clicca sull'intestazione di una colonna (es. stato, priorità, data creazione) per applicare un ordinamento.
\item Il sistema riordina la tabella in base ai valori della colonna selezionata e mostra un indicatore visivo della direzione dell'ordinamento (es. \uparrow o \downarrow).
\item \textit{[Opzionale]} L'utente clicca su ulteriori intestazioni di colonna per definire un ordinamento secondario o terziario.
\item \textit{[Opzionale]} Se l'utente clicca nuovamente su una colonna già ordinata, il sistema inverte la direzione dell'ordinamento o ripristina l'ordine predefinito.
\item Il sistema mostra la tabella aggiornata e, se previsto, memorizza i criteri di ordinamento nella sessione o nell'URL.
\end{enumerate}
\end{minipage} \\ \hline

\textbf{Estensioni} &
\begin{minipage}[t]{\linewidth}\raggedright
\textbf{3a. Ordinamento crescente/decrescente}\\
3a1. L'utente clicca una seconda volta sulla stessa intestazione di colonna.\\
3a2. Il sistema inverte la direzione dell'ordinamento (da crescente a decrescente o viceversa).\\
3a3. Il sistema aggiorna l'indicatore visivo (freccia \uparrow/\downarrow \) ) e ricarica la tabella.\\
3a4. Il flusso ritorna al termine del punto 3 dello scenario di successo.\\[0.6em]

\textbf{4a. Ordinamento secondario o multiplo}\\
4a1. L'utente clicca su un'altra intestazione di colonna diversa da quella attiva.\\
4a2. Il sistema imposta un secondo livello di ordinamento, mantenendo la gerarchia (es. priority → stato).\\
4a3. Il sistema aggiorna la tabella e mostra indicatori visivi multipli per le colonne coinvolte.\\
4a4. Il flusso ritorna al termine del punto 4 dello scenario di successo.\\[0.6em]

\textbf{4b. Rimozione di un ordinamento attivo}\\
4b1. L'utente clicca nuovamente su una colonna ordinata fino a tornare alla condizione “nessun ordine”.\\
4b2. Il sistema rimuove l'indicatore visivo e ripristina l'ordinamento predefinito.\\
4b3. Il flusso ritorna al termine del punto 3 dello scenario di successo.\\[0.6em]

\textbf{6a. Persistenza dell'ordinamento}\\
6a1. Il sistema salva automaticamente le informazioni sull'ordinamento corrente nella sessione o nell'URL.\\
6a2. Quando la pagina viene ricaricata o condivisa, l'ordinamento precedente viene ripristinato.\\
6a3. Il flusso prosegue dal punto 6 dello scenario di successo.\\[0.6em]

\textbf{6b. Ordinamento su colonna non ordinabile}\\
6b1. L'utente clicca su una colonna che non supporta l'ordinamento (es. colonna tag).\\
6b2. Il sistema ignora il click senza modificare la tabella.\\
6b3. Il flusso ritorna al punto 2 dello scenario di successo.
\end{minipage} \\ \hline

\textbf{Priorità} & Media \\ \hline
\textbf{Frequenza d'uso} & Tutti gli utenti, tra 1 e 3 volte al giorno per utente \\ \hline
\textbf{Informazioni associate} & Mostrato nella Tabella~\ref{tab:uc3-proprieta} \\ \hline

\end{longtable}
% Compilazione: XeLaTeX o LuaLaTeX consigliati (oppure pdfLaTeX con \usepackage[utf8]{inputenc})
% Preambolo richiesto:
% \usepackage{longtable}
% \usepackage{array}
% \usepackage{enumitem} % per itemize/enumerate compatti
\renewcommand{\arraystretch}{1.25}
\setlength{\tabcolsep}{8pt}

\begin{longtable}{|p{4cm}|p{11cm}|}
\caption{UC-3 Ordinamenti -- Template di Cockburn}\label{tab:uc3-cockburn}\\
\hline
\textbf{Campo} & \textbf{Valore} \\ \hline
\endfirsthead

\hline
\textbf{Campo} & \textbf{Valore} \\ \hline
\endhead
\hline
\multicolumn{2}{r}{\small Continua nella pagina successiva} \\
\endfoot
\hline
\endlastfoot

\textbf{ID - Nome dello UC} & 3 - Ordinamento \cite{uc3-figma} \\ \hline
\textbf{Data di creazione} & 14/10/2025 \\ \hline
\textbf{Attore primario} & Utente \\ \hline
\textbf{Trigger} & L'utente indica di voler ordinare l'elenco delle issues. \\ \hline
\textbf{Descrizione} & L'utente vuole ordinare le issues presenti nell'elenco tramite un criterio specifico. \\ \hline

\textbf{Precondizioni} &
\begin{minipage}[t]{\linewidth}\raggedright
\begin{itemize}[leftmargin=*,nosep]
\item PRE-1. L'utente è stato registrato da un amministratore ed ha effettuato correttamente l'autenticazione.
\item PRE-2. Il sistema ha caricato la pagina di riepilogo delle issues.
\end{itemize}
\end{minipage} \\ \hline

\textbf{Postcondizioni} & L'elenco delle issues è riordinato. \\ \hline

\textbf{Scenario di successo} &
\begin{minipage}[t]{\linewidth}\raggedright
\begin{enumerate}[leftmargin=*,nosep]
\item L'utente visualizza la tabella contenente le issues con le relative colonne.
\item L'utente clicca sull'intestazione di una colonna (es. stato, priorità, data creazione) per applicare un ordinamento.
\item Il sistema riordina la tabella in base ai valori della colonna selezionata e mostra un indicatore visivo della direzione dell'ordinamento (es. \uparrow o \downarrow).
\item \textit{[Opzionale]} L'utente clicca su ulteriori intestazioni di colonna per definire un ordinamento secondario o terziario.
\item \textit{[Opzionale]} Se l'utente clicca nuovamente su una colonna già ordinata, il sistema inverte la direzione dell'ordinamento o ripristina l'ordine predefinito.
\item Il sistema mostra la tabella aggiornata e, se previsto, memorizza i criteri di ordinamento nella sessione o nell'URL.
\end{enumerate}
\end{minipage} \\ \hline

\textbf{Estensioni} &
\begin{minipage}[t]{\linewidth}\raggedright
\textbf{3a. Ordinamento crescente/decrescente}\\
3a1. L'utente clicca una seconda volta sulla stessa intestazione di colonna.\\
3a2. Il sistema inverte la direzione dell'ordinamento (da crescente a decrescente o viceversa).\\
3a3. Il sistema aggiorna l'indicatore visivo (freccia \uparrow/\downarrow \) ) e ricarica la tabella.\\
3a4. Il flusso ritorna al termine del punto 3 dello scenario di successo.\\[0.6em]

\textbf{4a. Ordinamento secondario o multiplo}\\
4a1. L'utente clicca su un'altra intestazione di colonna diversa da quella attiva.\\
4a2. Il sistema imposta un secondo livello di ordinamento, mantenendo la gerarchia (es. priority → stato).\\
4a3. Il sistema aggiorna la tabella e mostra indicatori visivi multipli per le colonne coinvolte.\\
4a4. Il flusso ritorna al termine del punto 4 dello scenario di successo.\\[0.6em]

\textbf{4b. Rimozione di un ordinamento attivo}\\
4b1. L'utente clicca nuovamente su una colonna ordinata fino a tornare alla condizione “nessun ordine”.\\
4b2. Il sistema rimuove l'indicatore visivo e ripristina l'ordinamento predefinito.\\
4b3. Il flusso ritorna al termine del punto 3 dello scenario di successo.\\[0.6em]

\textbf{6a. Persistenza dell'ordinamento}\\
6a1. Il sistema salva automaticamente le informazioni sull'ordinamento corrente nella sessione o nell'URL.\\
6a2. Quando la pagina viene ricaricata o condivisa, l'ordinamento precedente viene ripristinato.\\
6a3. Il flusso prosegue dal punto 6 dello scenario di successo.\\[0.6em]

\textbf{6b. Ordinamento su colonna non ordinabile}\\
6b1. L'utente clicca su una colonna che non supporta l'ordinamento (es. colonna tag).\\
6b2. Il sistema ignora il click senza modificare la tabella.\\
6b3. Il flusso ritorna al punto 2 dello scenario di successo.
\end{minipage} \\ \hline

\textbf{Priorità} & Media \\ \hline
\textbf{Frequenza d'uso} & Tutti gli utenti, tra 1 e 3 volte al giorno per utente \\ \hline
\textbf{Informazioni associate} & Mostrato nella Tabella~\ref{tab:uc3-proprieta} \\ \hline

\end{longtable}
% Compilazione: XeLaTeX o LuaLaTeX consigliati (oppure pdfLaTeX con \usepackage[utf8]{inputenc})
% Preambolo richiesto:
% \usepackage{longtable}
% \usepackage{array}
% \usepackage{enumitem} % per itemize/enumerate compatti
\renewcommand{\arraystretch}{1.25}
\setlength{\tabcolsep}{8pt}

\begin{longtable}{|p{4cm}|p{11cm}|}
\caption{UC-3 Ordinamenti -- Template di Cockburn}\label{tab:uc3-cockburn}\\
\hline
\textbf{Campo} & \textbf{Valore} \\ \hline
\endfirsthead

\hline
\textbf{Campo} & \textbf{Valore} \\ \hline
\endhead
\hline
\multicolumn{2}{r}{\small Continua nella pagina successiva} \\
\endfoot
\hline
\endlastfoot

\textbf{ID - Nome dello UC} & 3 - Ordinamento \cite{uc3-figma} \\ \hline
\textbf{Data di creazione} & 14/10/2025 \\ \hline
\textbf{Attore primario} & Utente \\ \hline
\textbf{Trigger} & L'utente indica di voler ordinare l'elenco delle issues. \\ \hline
\textbf{Descrizione} & L'utente vuole ordinare le issues presenti nell'elenco tramite un criterio specifico. \\ \hline

\textbf{Precondizioni} &
\begin{minipage}[t]{\linewidth}\raggedright
\begin{itemize}[leftmargin=*,nosep]
\item PRE-1. L'utente è stato registrato da un amministratore ed ha effettuato correttamente l'autenticazione.
\item PRE-2. Il sistema ha caricato la pagina di riepilogo delle issues.
\end{itemize}
\end{minipage} \\ \hline

\textbf{Postcondizioni} & L'elenco delle issues è riordinato. \\ \hline

\textbf{Scenario di successo} &
\begin{minipage}[t]{\linewidth}\raggedright
\begin{enumerate}[leftmargin=*,nosep]
\item L'utente visualizza la tabella contenente le issues con le relative colonne.
\item L'utente clicca sull'intestazione di una colonna (es. stato, priorità, data creazione) per applicare un ordinamento.
\item Il sistema riordina la tabella in base ai valori della colonna selezionata e mostra un indicatore visivo della direzione dell'ordinamento (es. \uparrow o \downarrow).
\item \textit{[Opzionale]} L'utente clicca su ulteriori intestazioni di colonna per definire un ordinamento secondario o terziario.
\item \textit{[Opzionale]} Se l'utente clicca nuovamente su una colonna già ordinata, il sistema inverte la direzione dell'ordinamento o ripristina l'ordine predefinito.
\item Il sistema mostra la tabella aggiornata e, se previsto, memorizza i criteri di ordinamento nella sessione o nell'URL.
\end{enumerate}
\end{minipage} \\ \hline

\textbf{Estensioni} &
\begin{minipage}[t]{\linewidth}\raggedright
\textbf{3a. Ordinamento crescente/decrescente}\\
3a1. L'utente clicca una seconda volta sulla stessa intestazione di colonna.\\
3a2. Il sistema inverte la direzione dell'ordinamento (da crescente a decrescente o viceversa).\\
3a3. Il sistema aggiorna l'indicatore visivo (freccia \uparrow/\downarrow \) ) e ricarica la tabella.\\
3a4. Il flusso ritorna al termine del punto 3 dello scenario di successo.\\[0.6em]

\textbf{4a. Ordinamento secondario o multiplo}\\
4a1. L'utente clicca su un'altra intestazione di colonna diversa da quella attiva.\\
4a2. Il sistema imposta un secondo livello di ordinamento, mantenendo la gerarchia (es. priority → stato).\\
4a3. Il sistema aggiorna la tabella e mostra indicatori visivi multipli per le colonne coinvolte.\\
4a4. Il flusso ritorna al termine del punto 4 dello scenario di successo.\\[0.6em]

\textbf{4b. Rimozione di un ordinamento attivo}\\
4b1. L'utente clicca nuovamente su una colonna ordinata fino a tornare alla condizione “nessun ordine”.\\
4b2. Il sistema rimuove l'indicatore visivo e ripristina l'ordinamento predefinito.\\
4b3. Il flusso ritorna al termine del punto 3 dello scenario di successo.\\[0.6em]

\textbf{6a. Persistenza dell'ordinamento}\\
6a1. Il sistema salva automaticamente le informazioni sull'ordinamento corrente nella sessione o nell'URL.\\
6a2. Quando la pagina viene ricaricata o condivisa, l'ordinamento precedente viene ripristinato.\\
6a3. Il flusso prosegue dal punto 6 dello scenario di successo.\\[0.6em]

\textbf{6b. Ordinamento su colonna non ordinabile}\\
6b1. L'utente clicca su una colonna che non supporta l'ordinamento (es. colonna tag).\\
6b2. Il sistema ignora il click senza modificare la tabella.\\
6b3. Il flusso ritorna al punto 2 dello scenario di successo.
\end{minipage} \\ \hline

\textbf{Priorità} & Media \\ \hline
\textbf{Frequenza d'uso} & Tutti gli utenti, tra 1 e 3 volte al giorno per utente \\ \hline
\textbf{Informazioni associate} & Mostrato nella Tabella~\ref{tab:uc3-proprieta} \\ \hline

\end{longtable}
% Compilazione: XeLaTeX o LuaLaTeX consigliati (oppure pdfLaTeX con \usepackage[utf8]{inputenc})
% Preambolo richiesto:
% \usepackage{longtable}
% \usepackage{array}
% \usepackage{enumitem} % per itemize/enumerate compatti
\renewcommand{\arraystretch}{1.25}
\setlength{\tabcolsep}{8pt}

\begin{longtable}{|p{4cm}|p{11cm}|}
\caption{UC-3 Ordinamenti -- Template di Cockburn}\label{tab:uc3-cockburn}\\
\hline
\textbf{Campo} & \textbf{Valore} \\ \hline
\endfirsthead

\hline
\textbf{Campo} & \textbf{Valore} \\ \hline
\endhead
\hline
\multicolumn{2}{r}{\small Continua nella pagina successiva} \\
\endfoot
\hline
\endlastfoot

\textbf{ID - Nome dello UC} & 3 - Ordinamento \cite{uc3-figma} \\ \hline
\textbf{Data di creazione} & 14/10/2025 \\ \hline
\textbf{Attore primario} & Utente \\ \hline
\textbf{Trigger} & L'utente indica di voler ordinare l'elenco delle issues. \\ \hline
\textbf{Descrizione} & L'utente vuole ordinare le issues presenti nell'elenco tramite un criterio specifico. \\ \hline

\textbf{Precondizioni} &
\begin{minipage}[t]{\linewidth}\raggedright
\begin{itemize}[leftmargin=*,nosep]
\item PRE-1. L'utente è stato registrato da un amministratore ed ha effettuato correttamente l'autenticazione.
\item PRE-2. Il sistema ha caricato la pagina di riepilogo delle issues.
\end{itemize}
\end{minipage} \\ \hline

\textbf{Postcondizioni} & L'elenco delle issues è riordinato. \\ \hline

\textbf{Scenario di successo} &
\begin{minipage}[t]{\linewidth}\raggedright
\begin{enumerate}[leftmargin=*,nosep]
\item L'utente visualizza la tabella contenente le issues con le relative colonne.
\item L'utente clicca sull'intestazione di una colonna (es. stato, priorità, data creazione) per applicare un ordinamento.
\item Il sistema riordina la tabella in base ai valori della colonna selezionata e mostra un indicatore visivo della direzione dell'ordinamento (es. \uparrow o \downarrow).
\item \textit{[Opzionale]} L'utente clicca su ulteriori intestazioni di colonna per definire un ordinamento secondario o terziario.
\item \textit{[Opzionale]} Se l'utente clicca nuovamente su una colonna già ordinata, il sistema inverte la direzione dell'ordinamento o ripristina l'ordine predefinito.
\item Il sistema mostra la tabella aggiornata e, se previsto, memorizza i criteri di ordinamento nella sessione o nell'URL.
\end{enumerate}
\end{minipage} \\ \hline

\textbf{Estensioni} &
\begin{minipage}[t]{\linewidth}\raggedright
\textbf{3a. Ordinamento crescente/decrescente}\\
3a1. L'utente clicca una seconda volta sulla stessa intestazione di colonna.\\
3a2. Il sistema inverte la direzione dell'ordinamento (da crescente a decrescente o viceversa).\\
3a3. Il sistema aggiorna l'indicatore visivo (freccia \uparrow/\downarrow \) ) e ricarica la tabella.\\
3a4. Il flusso ritorna al termine del punto 3 dello scenario di successo.\\[0.6em]

\textbf{4a. Ordinamento secondario o multiplo}\\
4a1. L'utente clicca su un'altra intestazione di colonna diversa da quella attiva.\\
4a2. Il sistema imposta un secondo livello di ordinamento, mantenendo la gerarchia (es. priority → stato).\\
4a3. Il sistema aggiorna la tabella e mostra indicatori visivi multipli per le colonne coinvolte.\\
4a4. Il flusso ritorna al termine del punto 4 dello scenario di successo.\\[0.6em]

\textbf{4b. Rimozione di un ordinamento attivo}\\
4b1. L'utente clicca nuovamente su una colonna ordinata fino a tornare alla condizione “nessun ordine”.\\
4b2. Il sistema rimuove l'indicatore visivo e ripristina l'ordinamento predefinito.\\
4b3. Il flusso ritorna al termine del punto 3 dello scenario di successo.\\[0.6em]

\textbf{6a. Persistenza dell'ordinamento}\\
6a1. Il sistema salva automaticamente le informazioni sull'ordinamento corrente nella sessione o nell'URL.\\
6a2. Quando la pagina viene ricaricata o condivisa, l'ordinamento precedente viene ripristinato.\\
6a3. Il flusso prosegue dal punto 6 dello scenario di successo.\\[0.6em]

\textbf{6b. Ordinamento su colonna non ordinabile}\\
6b1. L'utente clicca su una colonna che non supporta l'ordinamento (es. colonna tag).\\
6b2. Il sistema ignora il click senza modificare la tabella.\\
6b3. Il flusso ritorna al punto 2 dello scenario di successo.
\end{minipage} \\ \hline

\textbf{Priorità} & Media \\ \hline
\textbf{Frequenza d'uso} & Tutti gli utenti, tra 1 e 3 volte al giorno per utente \\ \hline
\textbf{Informazioni associate} & Mostrato nella Tabella~\ref{tab:uc3-proprieta} \\ \hline

\end{longtable}